\section{Scheduled Actions (\texttt{ir.cron})}

\slidebackground{backgrounds/wallpaper.jpg}

\begin{frame}[fragile]{What Is \texttt{ir.cron}?}
	\begin{itemize}
		\item \textbf{\texttt{ir.cron}} is Odoo's scheduled action (cron job) model.
		\item It runs a Python method automatically on a schedule.
		\item Typical uses:
		\begin{itemize}
			\item send reminder emails
			\item archive old records
			\item compute nightly statistics
			\item sync data with external systems
		\end{itemize}
\begin{block}{Core Idea}
	A cron points to a model method and a timing interval.
\end{block}
	\end{itemize}
\end{frame}

\begin{frame}[fragile]{Defining a Cron in XML}
\begin{lstlisting}[style=xml]
<record id="ir_cron_student_reminder" model="ir.cron">
  <field name="name">Student: Send Attendance Reminder</field>
  <field name="model_id" ref="model_student_student"/>
  <field name="state">code</field>
  <field name="code">model._cron_send_attendance_reminder()</field>
  <field name="interval_number">1</field>
  <field name="interval_type">days</field>
  <field name="numbercall">-1</field>
  <field name="active" eval="True"/>
</record>
\end{lstlisting}
\end{frame}

\begin{frame}[fragile]{Breaking Down Cron Fields}
	\begin{itemize}
		\item \textbf{\texttt{name}}: label in Settings $\rightarrow$ Technical $\rightarrow$ Scheduled Actions
		\item \textbf{\texttt{model\_id}}: model that owns the method
		\item \textbf{\texttt{state}}: usually \texttt{code}
		\item \textbf{\texttt{code}}: Python expression executed by the cron
		\item \textbf{\texttt{interval\_number}} / \textbf{\texttt{interval\_type}}: how often it runs
		\item \textbf{\texttt{numbercall}}: how many times (\texttt{-1} = forever)
		\item \textbf{\texttt{active}}: enable/disable the cron
	\end{itemize}
\end{frame}

\begin{frame}[fragile]{Interval Types}
	\begin{itemize}
		\item Common \texttt{interval\_type} values:
		\begin{itemize}
			\item \texttt{minutes}
			\item \texttt{hours}
			\item \texttt{days}
			\item \texttt{weeks}
			\item \texttt{months}
		\end{itemize}
	\end{itemize}
\begin{lstlisting}[style=xml]
<field name="interval_number">30</field>
<field name="interval_type">minutes</field>
\end{lstlisting}
	\begin{itemize}
		\item Choose intervals carefully -- too frequent jobs can overload the server.
	\end{itemize}
\end{frame}

\begin{frame}[fragile]{Writing the Cron Method}
\begin{block}{Method Syntax}
\begin{lstlisting}
@api.model
def _cron_send_attendance_reminder(self):
	students = self.search([
		('active', '=', True),
		('email', '!=', False),
	])
	template = self.env.ref(
		'school_manager.mail_template_attendance_reminder',
		raise_if_not_found=False,
	)
	if not template:
		return False
	for student in students:
		template.send_mail(student.id, force_send=False)
	return True
\end{lstlisting}
\end{block}
\end{frame}

\begin{frame}[fragile]{Important Cron Method Rules}
	\begin{itemize}
		\item Use \texttt{@api.model} for cron entry methods.
		\item Keep the method name private-style: \texttt{\_cron\_...}
		\item Make the job \textbf{idempotent} when possible (safe if run twice).
		\item Avoid huge work in one transaction; batch when needed.
		\item Log failures clearly.
	\end{itemize}
\begin{lstlisting}
@api.model
def _cron_archive_old_logs(self):
	old = self.env['student.log'].search([
		('create_date', '<', fields.Datetime.subtract(
			fields.Datetime.now(), days=90)),
	])
	old.unlink()
\end{lstlisting}
\end{frame}

\begin{frame}[fragile]{Manual Testing of a Cron}
\begin{lstlisting}
# Odoo shell
env['student.student']._cron_send_attendance_reminder()

# Or open Settings > Technical > Automation > Scheduled Actions
# and click Run Manually
\end{lstlisting}
	\begin{itemize}
		\item Always test the method manually before relying on the schedule.
		\item Check server logs after the first automatic runs.
	\end{itemize}
\end{frame}

\begin{frame}[fragile]{User, Company, and Rights in Crons}
	\begin{itemize}
		\item Cron jobs usually run as the OdooBot / superuser environment.
		\item Still write safe domains and business checks.
		\item If you need a specific company:
\begin{lstlisting}
@api.model
def _cron_company_cleanup(self):
	for company in self.env['res.company'].search([]):
		students = self.with_company(company).search([
			('company_id', '=', company.id),
			('active', '=', False),
		])
		# company-specific processing
\end{lstlisting}
	\end{itemize}
\end{frame}

\begin{frame}{Cron Best Practices}
	\begin{itemize}
		\item One cron = one clear responsibility.
		\item Prefer nightly/hourly jobs over very short loops.
		\item Use \texttt{numbercall = -1} for recurring jobs.
		\item Keep XML cron definitions in \texttt{data/} with \texttt{noupdate} when users may customize schedules.
		\item Monitor runtime and failure logs in production.
	\end{itemize}
\end{frame}
